\documentclass{article}

% Language setting
% Replace `english' with e.g. `spanish' to change the document language
\usepackage[english]{babel}

% Set page size and margins
% Replace `letterpaper' with `a4paper' for UK/EU standard size
\usepackage[a4paper,top=2cm,bottom=2cm,left=3cm,right=3cm,marginparwidth=1.75cm]{geometry}

% Useful packages
\usepackage{amsmath}
\usepackage{graphicx}
\usepackage[colorlinks=true, allcolors=blue]{hyperref}

\title{Discrete Math Course}
\author{Wisam Yassar Mahardika}
\date{\today}

\begin{document}
\maketitle
\clearpage
\tableofcontents
\clearpage
\section{Propositions, Negations, Conjunctions and Disjunctions}
\subsection{Propositions}
A \textbf{Proposition} is a declarative statement that is either true or false.
\begin{itemize}
\item The sky is blue
\item The moon is made of cheese
\item Luke, I am your father
\item Sit down
\item $ X + 1 = 2 $ \\
\end{itemize}

 \textit{Sit down is neither true or false so it is not a proposition text} \\
\textit{$ X + 1 = 2 $ is also neither true or false because we don't know the value of x}
\\

\subsection{Constructing Compound Propositions}
A \textbf{Compound Proposition} is comprised of propositions and one or more of the following connectives:
\\


\begin{itemize}
\item Negation $ \neg $ ``not`
\item Conjunction $ \land $     ``and``
\item Disjunction $ \lor $ ``or``
\item Implication $ \implies $ ``if, then``
\item Biconditional $ \Longleftrightarrow{} $ ``if and only if``\\
\end{itemize}

Each proposition is represented by a propositional variable \textit{p, q, r, s ...}.\\


Example: $ p \implies q $

\subsubsection{Negation}
The \textbf{negation} of the proposition $ p $ is $ \neg p $ (not $p$) \\


Example if $ p $ denotes ``The grass is green``, then $ \neg p $ denotes ``it is not the case that the grass is green`` or more simply ``the grass is not green``.\\


\begin{itemize}
\item My dog is the cutest dog  ($ \neg p $ My dog is \textbf{not} the cutest dog)
\item The door is not open ($ \neg p $ The door \textbf{is} open)
\item Are we there yet
\end{itemize}
\textit{The third statement is not a proposition because it is nor true or false.}

\clearpage

\subsubsection{Truth Table}
Each row of a \textbf{truth table} gives us one possibility for the truth values of our proposition(s), since each proposition has two possible truth value, \textit{true} or \textit{false}, we will have 2 rows for each propositions (or $ 2^n $ rows where $ n = \# $ of propositions). \\


\textit{Truth Table for $ \neg p $} 
\begin{tabular}{|c|c|}
\hline
$ p $ & $ \neg p $ \\
\hline
T & F \\
F & T \\
\hline
\end{tabular} \\
\textit{$ p $ is combinations and $ \neg p $ is connectives} \\


\subsubsection{Conjunction}
The \textbf{conjunction} of propositions $ p $ and $ q $ is denoted  $ p \land q $ and read $ p $ ``and`` $ q $.
For a conjunction to be true, \textbf{both} propositions must be true. \\

\textit{Truth Table for $ p \land q $} 
\begin{tabular}{|c|c|c|}
\hline
$ p $ & $ q $ & $ p \land q $ \\
\hline
T & T & T \\
T & F & F \\
F & T & F \\
F & F & F \\
\hline
\end{tabular}\\


\subsubsection{Disjunction}
The \textbf{disjunction} of propositions $ p $ and $ q $ is denoted $ p \lor q $ and read $ p $ ``or`` $ q $.
For a disjunction to be true. \textbf{either} proposition must be true.\\


\textit{Truth Table for $ p \lor q $}
\begin{tabular}{|c|c|c|}
\hline
$ p $ & $ q $ & $ p \lor q $ \\
\hline
T & T & T \\ 
T & F & T \\
F & T & T \\
F & F & F \\
\hline
\end{tabular}\\


\subsubsection{The Connective ``or`` in English ``xor``}
\begin{itemize}
	\item \textbf{Inclusive ``or`` $ p \lor q $}
	The prerequisite for MA420 is either MA315 or MA335.
	\item \textbf{Exclusive ``or`` $ p \oplus q $}
	You get soup or salad with your entree.
\end{itemize}


\textit{Truth Table for $ p \oplus q $}
\begin{tabular}{|c|c|c|}
\hline
$ p $ & $ q $ & $ p \oplus q $ \\
\hline
T & T & F \\ 
T & F & T \\
F & T & T \\
F & F & F \\
\hline
\end{tabular}\\

\clearpage
\section{Implications (converse, inverse, contrapositive) and biconditionals}
\subsection{Implication (conditional statements)}
The \textbf{implication} of propositions $ p $ and $ q $ is denoted $ p \implies q $ and read ``if $ p $ then $ q $`` or ``$ p $ implies $ q $``.\\


When the hypothesis is true, the conclusion must be true for the implication to be true. when the hypothesis is false, the conclusion is true.\\


$ p $ denotes ``It is a holiday`` $ q $ denotes ``The store is closed``. \\


\textit{Truth Table for $ p \implies q $}
\begin{tabular}{|c|c|c|}
\hline
$ p $ & $ q $ & $ p \implies q $ \\
\hline
T & T & T \\ 
T & F & F \\
F & T & T \\
F & F & T \\
\hline 
\end{tabular}
\subsubsection{Converse, Inverse, Contrapositive}
From our implication, $ p \implies q $, we can construct 3 new conditional statements.
\begin{enumerate}
	\item \textbf{Converse: }$ q \implies p $ 
	\item \textbf{Inverse: }$ \neg p \implies \neg q $
	\item \textbf{Contrapositive: }$ \neg q \implies \neg p $
\end{enumerate}


\textit{e.g} It is raining is a sufficient condition for me not going to town.
\begin{itemize}
	\item in an ``if then`` form it looks like this: \textbf{If} it is raining, \textbf{then} I won't go to town. $ p \implies q $
	\item In converse it looks like this: \textbf{If} I don't go to town, \textbf{then} it is raining.
	\item In inverse it looks like this: \textbf{If} it is not raining, \textbf{then} I will go to town.
	\item In contrapositive it looks like this: \textbf{If} I will go to town, \textbf{then} it is not raining.\\
\textbf{Note: }Contrapositive always have the same \textit{truth value} as $ p \implies q $.
\end{itemize}
\subsection{Biconditional}
The \textbf{bicondtional} of proposition $ p $ and $ q $ is denoted $ p \Longleftrightarrow{} q $ and read ``$ p $ if and only if $ q $``.\\


For a biconditional to be true, both propositions must share the same truth value. \\


\textit{Truth Table for $ p \Longleftrightarrow q $}
\begin{tabular}{|c|c|c|}
\hline
$ p $ & $ q $ & $ p \Longleftrightarrow q $ \\
\hline
T & T & T \\ 
T & F & F \\
F & T & F \\
F & F & T \\
\hline 
\end{tabular}


Our biconditional $ p \Longleftrightarrow q $ can also be written as a compound proposition. \\
$ (p \Longleftrightarrow q) \equiv (p \implies q) \land (q \implies p) $ \\


\begin{tabular}{|c|c|c|c|c|c|}
\hline
$ p $ & $ q $ & $ p \implies q $ & $ q \implies p $ & $ (p \implies q) \land (q \implies p) $ & $ p \Longleftrightarrow q $ \\
\hline
T & T & T & T & T & T \\ 
T & F & F & T & F & F \\
F & T & T & F & F & F \\
F & F & T & T & T & T \\
\hline 
\end{tabular}

\clearpage
\section{Constructing A Truth Table for Compound Propositions}
\subsection{Compound Proposition Truth Table Walk Thru}
\begin{itemize}
	\item \textbf{Rows}
		\begin{itemize}
			\item Need a row for every possible combination of value for the compound propositions.
		\end{itemize}
	\item \textbf{Columns}
		\begin{itemize}
			\item Need a column for each propositional variable
			\item Need a column for the truth value of each expression that occurs in the compound proposition as it is built up
			\item Need a column for the compound proposition (usually at far right)
		\end{itemize}		
	\item \textbf{Order Of Operations} \\
		\begin{tabular}{|c|c|}
			\hline
			Precedence & Operator\\
			\hline
			1 & $ \neg $ \\
			2 & $ \land $ \\
			3 & $ \lor $ \\
			4 & $ \implies $ \\
			5 & $ \Longleftrightarrow $ \\
			\hline
		\end{tabular}
\end{itemize}
\subsubsection{Construct a table for $ p \lor q \implies \neg r $}
\begin{enumerate}
	\item First, construct columns for each proposition; $ p,q,r $ \textit{these may be referred as atomic propositions} \\
		

		\begin{tabular}{|c|c|c|}
			\hline
			$ p $ & $ q $ & $ r $ \\
			\hline
		\end{tabular}
	\item Next, create a column for each compound proposition; $ p \lor q \implies \neg r $ \\
		

		\begin{tabular}{|c|c|c|c|c|}
			\hline
			$ p $ & $ q $ & $ r $ & $ p \lor q $ & $ \neg r $ \\
			\hline
		\end{tabular}
	\item Lastly, create a column for the final compound proposition \\
		

		\begin{tabular}{|c|c|c|c|c|c|}
			\hline
			$ p $ & $ q $ & $ r $ & $ p \lor q $ & $ \neg r $ & $ p \lor q \implies \neg r $\\
			\hline
		\end{tabular}
	\item Now create as many rows as necessary for all possible combinations of expressions that occurs in your compound propositions, Because we have 3 propositions, we need $2^3 = 8$ rows. \\


		\begin{tabular}{|c|c|c|c|c|c|}
			\hline
			$ p $ & $ q $ & $ r $ & $ p \lor q $ & $ \neg r $ & $ p \lor q \implies \neg r $\\
			\hline
			$ $ & $ $ & $ $ & $ $ & $ $ & $ $\\
			$ $ & $ $ & $ $ & $ $ & $ $ & $ $\\
			$ $ & $ $ & $ $ & $ $ & $ $ & $ $\\
			$ $ & $ $ & $ $ & $ $ & $ $ & $ $\\
			$ $ & $ $ & $ $ & $ $ & $ $ & $ $\\
			$ $ & $ $ & $ $ & $ $ & $ $ & $ $\\
			$ $ & $ $ & $ $ & $ $ & $ $ & $ $\\
			$ $ & $ $ & $ $ & $ $ & $ $ & $ $\\
			\hline
		\end{tabular}
	\item Fill in the table \\


		\begin{tabular}{|c|c|c|c|c|c|}
			\hline
			$ p $ & $ q $ & $ r $ & $ p \lor q $ & $ \neg r $ & $ p \lor q \implies \neg r $\\
			\hline
			T & T & T & T & F & F\\
			T & T & F & T & T & T\\
			T & F & T & T & F & F\\
			T & F & F & T & T & T\\
			F & T & T & T & F & F\\
			F & T & F & T & T & T\\
			F & F & T & F & F & T\\
			F & F & F & F & T & T\\
			\hline
		\end{tabular}
\end{enumerate}
\subsection{Practice}
Create a truth table for $(p \lor \neg q) \implies (p \land q)$ \\


\textbf{Solution:}
we have 2 variables which is p and q, so $ n = 2 $\\
$2^2 = 4$\\
So we need 4 rows \\


\begin{tabular}{|c|c|c|c|c|c|}
	\hline
	$ p $ & $ q $ & $ \neg q $ & $ p \lor \neg q $ & $ p \land q $ & $(p \lor \neg q) \implies (p \land q)$ \\
	\hline
	T & T & F & T & T & T \\
	T & F & T & T & F & F \\
	F & T & F & F & F & T \\
	F & F & T & T & F & F \\
	\hline
\end{tabular}
\end{document}
